\section{Постановка задачи}
\label{sec:Chapter1} \index{Chapter1}

%Необходимо формально изложить суть задачи в данной секции, предоставив такие
%ясные и точные описания, которые позволят в последующем оценить, насколько
%разработанное решение соответствует поставленной задаче. Текст главы должен
%следовать структуре технического задания, включая как описание самой задачи,
%так и набор требований к ее решению.

Обозначим $\mathcal{R}_q^n = \Z_q[x]/(x^n+1)$, где $q$ --- простое и $n$ --- степень $2$. Тогда многочлен деления круга $x^n+1$ неприводим над $\Z$, откуда $\Z[x]/(x^n+1)$ --- циклотомическое поле, а само $\mathcal{R}_q^n$ --- подходящее по структуре кольцо.

Будем обозначать биномиальное распределение c $n$ попытками и вероятностью успеха $p$, как $\mathcal{B}(n, p)$, а его значение в точке $x$ за $\mathcal{B}(n, p)[x]$ (здесь $0 \leq x \leq n$). Если $\xi, \mu \sim \mathcal{B}(k, p)$ --- независимые случайные величины, распределение разности $\xi-\mu$ обозначим за $\chi_k$ ({\it centered binomial distribution}), а его значение в точке $x$ за $\chi_k[x]$ (здесь $-k \leq x \leq k$).
